\documentclass[11pt, a4paper]{article}
\usepackage[margin=1in, top=1.2in, bottom=1.2in]{geometry}
\usepackage{amsmath}
% ---- Font & Micro-Typography ----
\usepackage[T1]{fontenc}
\usepackage{lmodern}
\usepackage{microtype}
% ---- Formatting & Aesthetics ----
\usepackage{graphicx}
% ---- Bibliography ----
\usepackage[backend=biber, style=numeric-comp, sorting=nyt, giveninits=true]{biblatex}


% ====================================================================================
\begin{document}

\begin{center}
  \LARGE\textbf{Daily Rant Material Strength and Simulation}\\
  Tadouanla Guetchuin Billy\\
  28177
\end{center}

\begin{itemize}
  \item \textbf{Introduction to Matrices:} How to compute the inverse of a matrix.
        \begin{align*}
          \mathbf{l_1}: y = a_1x + b_1 \rightarrow y - a_1x  = b_1 \\
          \mathbf{l_2}: y = a_2x + b_2  \rightarrow y - a_2x = b_2 \\
          \begin{bmatrix}
            -a_1 & 1 \\
            -a_2 & 1
          \end{bmatrix}
          \begin{pmatrix}
            x \\
            y
          \end{pmatrix}
          =
          \begin{bmatrix}
            b1 \\
            b2
          \end{bmatrix}
        \end{align*}
        \textbf{Inverse}
        $$
          \begin{bmatrix}
            A
          \end{bmatrix}
          \begin{bmatrix}
            A
          \end{bmatrix}^{-1} =
          \begin{bmatrix}
            A
          \end{bmatrix}^{-1}
          \begin{bmatrix}
            A
          \end{bmatrix} =
          \begin{bmatrix}
            I
          \end{bmatrix}
        $$

\end{itemize}
We worked on forces namely internal forces on a bar and static equilibrium. Below is a diagram
showing the different loads/forces acting on a bar at fixed on both ends. We used 3 different
approaches in solving for the \textbf{stress} and \textbf{deformation}. Stress and deformation
are given resolved using the formulas below:
\begin{align*}
  \sigma & = \frac{N}{A}\hspace{2cm} \text{Stress}                           \\
  A      & = \frac{\pi d^2}{4} \hspace{1.7cm} \text{Cross sectional Area}    \\
  \delta & = \sum\frac{F\cdot L}{E\cdot A} \hspace{1.0cm} \text{Deformation}
\end{align*}

\noindent Below is a diagram showing all the forces acting on compound bar.
%------ First Picture -------------
\begin{figure}[h!]
  \center
  \includegraphics[scale=0.35]{./figures/MSS_1.jpg}
\end{figure}

\noindent \textbf{Analytical:} The problem was approached analytically using the principles of
classical mechanics, specifically static and solid mechanics.\\
\noindent \textbf{Numerical:} The numerical method was developed to simulate the physical system
and solve the unknown stresses and deformations using matrices.\\
\noindent \textbf{Kraken Program:} Finally a specialised model built using Kraken(Program from
the professor himself). The program is capable of computing the resulting forces on the system
and even drawing the physical system itself.


%------ Second Picture -------------
% \begin{figure}[h!]
%   \center
%   \includegraphics[scale=0.5]{./figures/MSS_2.jpg}
% \end{figure}

\end{document}
